% Nouveau EGA V (cote 25): archive cote 25, site page 156 (PDF page 157).
% typescript. Draft compared: none.
% AI-transcribed by Claude Opus 5.5 (claude-opus-5-5, adaptive thinking, effort high), Message Batches API, batch ega25b (msgbatch_01A7n5YuXubUPNK1zsepi8RG), 2026-11-13; not yet checked by eye.
\scanpage{25/156}
\mspage{\hand{17}}

\xout{\unsure{g}} ``générique'')\hand{\unsure{d}}.   Ceci posé, il faudrait vérifier que l'ensem ble $V^2$
des points supersinguliers de u rel. à Y n'est autre que l'ensemble des
points en lesquels $V^{(1)} \hand{\longrightarrow} Y$ est ramifié, et préciser ce point
comme une identité de sous-préschémas fermés de X , $V^{\unsure{\over{2}{2}}}$ \over{=}{=} \hand{``}sous-préschéma
de ramification\hand{''} de $V^1 \hand{\longrightarrow} Y$. Je n'ai pas fait l'exercice en détail,
mais ne doute pas qu'il ne puisse se faire par essentiellement l\unsure{a|e} même
\xout{méthode} \inserted[typed]{genre de dévissage} que cel\unsure{le|ui} développé\xout{e} après la prop.7.

\hand{\unsure{d}} Si $E \hand{\longrightarrow} F$ est un homomorphisme de Modules localement libres \unsure{s|d}\unsure{u|a}r \inserted[typed]{le préschéma régulier}
X, de rangs m et d ($m \hand{\leqslant} d$), et si \unsure{Z|Y} est l'ensemble des points de X en
lesquels le rang de u est $\hand{\leqslant}$ m-1, alors \inserted[typed]{on peut montrer que} Z est de codimension $\hand{\geqslant}$ d-m+1
en chacun de ses points, \xout{comme on vérifie immédiatement par récurrence sur}
\xout{m, à partir du cas m=1 (}

\mspage{\hand{156}}
